\subsection*{MT-S506. Dirichlet support of possible poles for every $S_5$ extension}

Let $L/\mathbb Q$ be a Galois extension with group $G=S_5$,
let $E=L^{A_5}$, and let $\eta$ be its quadratic character.
No restriction on complex conjugation or on ramification at $5$
is imposed in this subsection.  Use the irreducible representations
\[
 \begin{gathered}
 \mathbf1,\quad\eta,\quad V,\quad V\eta,\quad W,\quad W\eta,\quad X,\\
 V=\mathbf C^5/\mathbf C\mathbf1,\quad
 W=\operatorname{Sym}^2V-\mathbf1-V,\quad X=\bigwedge^2V.
 \end{gathered}
 \tag{S506.1}
\]
The notation for $W$ first denotes a virtual representation;
the character calculation below proves that it is an actual
irreducible representation of dimension $5$.  These are
Calegari's $\rho_4,\rho_5,\rho_6$, respectively, and $X\eta=X$.
All $L$-functions here include the finite factors at ramified primes.

For $z$ in the open critical strip $0<\Re z<1$, put
\[
 (a,b,c,d,x,y,e)=\operatorname{ord}_{s=z}
 (\zeta_{\mathbb Q},L_\eta,L_V,L_{V\eta},L_W,L_{W\eta},L_X).
 \tag{S506.2}
\]
An order is positive at a zero and negative at a pole.
Unconditionally $L_X$ is an entire Hecke $L$-function of a
nontrivial finite-order character over a degree-six field.
Consequently $a,b,e\ge0$, and the pole orders satisfy
\[
 \begin{aligned}
 \max(0,-c)&\le\min(a,e),&
 \max(0,-d)&\le\min(b,e),\\
 \max(0,-x)&\le b,&
 \max(0,-y)&\le a.
 \end{aligned}
 \tag{S506.3}
\]
In particular every possible pole of an irreducible Artin factor
in this strip lies in
\[
          Z(\zeta_{\mathbb Q})\ \cup\ Z(L_\eta).
 \tag{S506.4}
\]
Under the designated input $A_D$, every such pole therefore lies
on $\Re z=1/2$.  This conclusion uses only rational Dirichlet GRH.
It neither asserts that these poles exist nor proves their absence,
and it gives no general zero-location theorem for the other factors.

\begin{proof}
The complete character table, in the order (S506.1), is
\[
\begin{array}{c|r|rrrrrrr}
\text{cycle type}&\text{class size}&\mathbf1&\eta&V&V\eta&W&W\eta&X\\\hline
1^5&1&1&1&4&4&5&5&6\\
2\,1^3&10&1&-1&2&-2&1&-1&0\\
2^2\,1&15&1&1&0&0&1&1&-2\\
3\,1^2&20&1&1&1&1&-1&-1&0\\
3\,2&20&1&-1&-1&1&1&-1&0\\
4\,1&30&1&-1&0&0&-1&1&0\\
5&24&1&1&-1&-1&0&0&1
\end{array}
 \tag{S506.5}
\]
Indeed $\chi_V(g)=\#\operatorname{Fix}(g)-1$ and
\[
 \chi_W(g)=\frac{\chi_V(g)^2+\chi_V(g^2)}2-1-\chi_V(g),\qquad
 \chi_X(g)=\frac{\chi_V(g)^2-\chi_V(g^2)}2.
\]
These formulas give every displayed entry.  Weighted inner products
of the seven columns are $\delta_{ij}$: multiplication by the
class-size diagonal matrix gives $120I$ for the Gram matrix.
A virtual character of norm one is the positive or negative of
an irreducible character; its positive value at the identity fixes
the positive sign.  Thus all seven columns are actual irreducibles.
Their squared dimensions sum to $120$, so the list is complete.
This also proves $X\eta=X$ directly, since $X$ vanishes on odd
permutations.

Let $J=S_4$ be a point stabilizer and let
$F_{20}=\operatorname{AGL}_1(\mathbb F_5)\subset S_5$ be the
normalizer of a $5$-cycle.  Write an affine element as
$t\mapsto ut+v$, and let $\tau(u)$ be a primitive quartic character
of $\mathbb F_5^\times$, inflated to $F_{20}$.  The sign restricted
to $F_{20}$ is $\tau^2$.  Let
$T=S_2\times S_3$ be a two-subset stabilizer and let
$\alpha$ be sign on its $S_2$ factor and trivial on its $S_3$ factor.
Frobenius reciprocity and (S506.5) give
\[
\begin{array}{c|l}
\text{induced character}&\text{decomposition in }G\\\hline
\operatorname{Ind}_{J}^{G}\mathbf1&\mathbf1+V\\
\operatorname{Ind}_{J}^{G}(\eta|_J)&\eta+V\eta\\
\operatorname{Ind}_{F_{20}}^{G}\mathbf1&\mathbf1+W\eta\\
\operatorname{Ind}_{F_{20}}^{G}(\eta|_{F_{20}})&\eta+W\\
\operatorname{Ind}_{F_{20}}^{G}\tau&X\\
\operatorname{Ind}_{T}^{G}\alpha&V+X\\
\operatorname{Ind}_{T}^{G}(\alpha\eta|_T)&V\eta+X.
\end{array}
 \tag{S506.6}
\]
Here is an explicit check of the potentially sensitive quartic
identity.  The affine group has the identity, four nonidentity
translations of cycle type $5$, five elements of type $2^2\,1$,
and ten of type $4\,1$.  The first two sets have multiplier $1$,
the third has multiplier $-1$, and the last two multiplier classes
have $\tau$-values $i$ and $-i$.  For $X$ the weighted inner-product
numerator is $6+4+10=20$.  For every other column the numerator
is zero, because the two quartic multiplier classes cancel and
the remaining two multiplier contributions cancel.  Division by
$|F_{20}|=20$ proves the identity.  The other rows follow by the
same finite averaging formula
\[
 \bigl\langle\operatorname{Ind}_{H}^{G}\lambda,R\bigr\rangle_G
       =\frac1{|H|}\sum_{h\in H}\overline{\lambda(h)}\chi_R(h).
 \tag{S506.7}
\]
For the two $T$ rows, $\alpha$ and $\alpha\eta$ are respectively
sign on the $S_2$ and on the $S_3$ factor, which fixes the twist
convention.

For a linear character $\lambda$ of any subgroup $H\subset G$,
Artin induction and class field theory identify
$L(s,\operatorname{Ind}_H^G\lambda)$ with its Hecke $L$-function
over $L^H$.  It is holomorphic in the open critical strip,
including when $\lambda$ is trivial, in which case it is a
Dedekind zeta function.  For $\tau\ne\mathbf1$ it is entire.
Taking orders in (S506.6) yields
\[
 \begin{gathered}
 a+c\ge0,\quad b+d\ge0,\quad a+y\ge0,\quad b+x\ge0,\\
 e\ge0,\quad c+e\ge0,\quad d+e\ge0.
 \end{gathered}
\]
The orders exist as integers because Brauer induction gives
meromorphic continuation of every Artin factor.  Since $a,b,e$
are nonnegative, these inequalities prove (S506.3)--(S506.4).
The last assertion is precisely the Dirichlet zero-location
input $A_D$ applied to the trivial and the quadratic characters.
\end{proof}

The classical all-prime induction and multiplication identities
are Prasad--Yogananda, \emph{A Report on Artin's holomorphy
conjecture}, Section~2.1, Properties~1--2, PDF p.~3;
Hecke entireness is Section~3.1, pp.~4--5; meromorphicity is
Theorem~3.3, p.~5.  The character labels agree with Calegari,
\emph{The Artin conjecture for some $S_5$-extensions},
author-hosted publisher PDF p.~4, Section~2 character table.
The character calculations here are proved explicitly; the
classical analytic inputs are cited, not re-proved.

\subsection*{MT-S507. The complete monomial cone and the remaining five-dimensional gate}

Retain (S506.1)--(S506.2), and assume additionally that
$L_V$ and $L_{V\eta}$ are holomorphic at $z$, that is, $c,d\ge0$.
This is only a pointwise hypothesis for the following cone theorem.
It holds at every point for the Calegari family of MT-S501:
complex conjugation has cycle type $(12)(34)$, $5$ is unramified,
and $\operatorname{Frob}_5$ has the same cycle type.
Calegari's Theorem~1.1, author-hosted publisher PDF p.~2, supplies
this four-dimensional holomorphy unconditionally.  No GRH input
is needed for the cone theorem or its numerical bounds.

At $z$, all tests
\[
 \operatorname{ord}_z L\bigl(s,\operatorname{Ind}_H^G\lambda\bigr)
             \ge0\qquad(H\subset G,\ \dim\lambda=1)
 \tag{S507.1}
\]
are equivalent, after imposing $c,d\ge0$, to the following twelve
linear inequalities:
\[
 \begin{gathered}
      a,b,c,d,e\ge0,\\
 g_1=x+y\ge0,\quad g_2=b+x\ge0,\quad g_3=a+y\ge0,\\
 g_4=c+x+e\ge0,\quad g_5=d+y+e\ge0,\\
 g_6=a+c+x\ge0,\quad g_7=b+d+y\ge0.
 \end{gathered}
 \tag{S507.2}
\]
The assertion is an exact description of this cone, not a claim
that all of its integer points occur as arithmetic order vectors.
In particular $c,d\ge0$ are extra holomorphy inputs, not deductions
from the monomial tests alone.

Define the relative and total Dedekind orders
\[
 q=\operatorname{ord}_z\frac{\zeta_L}{\zeta_E}
       =4(c+d)+5(x+y)+6e,
 \qquad
 r=\operatorname{ord}_z\zeta_L=a+b+q.
 \tag{S507.3}
\]
If $L_W$ has a pole of order $v>0$ at $z$, then
\[
 \begin{gathered}
 x=-v,\quad y\ge v,\quad b\ge v,\quad c+e\ge v,
       \quad a+c\ge v,\\
 v\le\min\{b,c+e,a+c\},\qquad q\ge4v,\qquad r\ge5v.
 \end{gathered}
 \tag{S507.4}
\]
For a pole of $L_{W\eta}$, interchange
$(a,b),(c,d),(x,y)$.  Both five-dimensional factors cannot have
a pole at the same point.  In particular $q\le3$ implies that
all irreducible factors are holomorphic at $z$, and every possible
pole order is at most $\lfloor q/4\rfloor$.
In the branch $c=0$ of (S507.4),
\[
        a,b,e,y\ge v,\qquad q\ge6v,\qquad r\ge8v.
 \tag{S507.5}
\]
The analogous assertion with $d=0$ holds for $W\eta$.
All these numerical constants are sharp for the exact cone
(S507.2), including the extra branch condition in (S507.5).

Let $Z(f)$ denote its zero set within the open critical strip.
For this holomorphy family the corresponding precise necessary
simultaneous-zero conditions are
\[
\begin{aligned}
 \operatorname{Pole}(L_W)&\subseteq
 Z(L_\eta)\cap\bigl(Z(L_V)\cup
             (Z(\zeta_{\mathbb Q})\cap Z(L_X))\bigr),\\
 \operatorname{Pole}(L_{W\eta})&\subseteq
 Z(\zeta_{\mathbb Q})\cap\bigl(Z(L_{V\eta})\cup
             (Z(L_\eta)\cap Z(L_X))\bigr).
\end{aligned}
 \tag{S507.6}
\]
Under $A_D$ these possible poles lie on the critical line by
MT-S506.  Under $M_H$ and the Calegari local hypotheses they
must additionally be nonreal, by MT-S505.  Neither hypothesis
by itself supplies the missing noncoincidence statement in
(S507.6), and no simplicity of any zero is asserted here.

\begin{proof}[Proof of the exact cone]
Besides the rows already established in (S506.6), use the
subgroup $A_4$ fixing the fifth letter and a nontrivial
character $\nu$ of its quotient $A_4/V_4=C_3$; the two-subset
stabilizer $T$; and
\[
  D_8=\langle r_0,s_0\rangle\subset S_4,\qquad
  r_0=(1234),\quad s_0=(13),\qquad
  \lambda(r_0)=-1,\quad\lambda(s_0)=1.
\]
These values respect $r_0^4=s_0^2=1$ and
$s_0r_0s_0=r_0^{-1}$, so they define a linear character.
Formula (S506.7) and the table give
\[
 \begin{gathered}
 \operatorname{Ind}_{A_4}^{G}\nu=W+W\eta,\qquad
 \operatorname{Ind}_{D_8}^{G}\lambda=V+W+X,\\
 \operatorname{Ind}_{T}^{G}\mathbf1=\mathbf1+V+W.
 \end{gathered}
 \tag{S507.7}
\]
Twisting by $\eta$ gives the other two needed rows, since
$X\eta=X$.  For the $A_4$ identity, its identity and three
involutions contribute $5+3$ to each rank-five inner product;
the eight $3$-cycles split into two sets of four with character
values $\omega,\omega^2$, so their weighted contribution is
$-4(\omega+\omega^2)=4$.  Division by $12$ gives one.
The other inner products vanish.  For the $D_8$ identity,
its elements have cycle types $1$, two $4$-cycles, three double
transpositions, and two transpositions.  Their $\lambda$-weights
sum, within those types, to $1,-2,-1,2$, respectively.
Against the seven columns this yields
$(0,0,1,0,1,0,1)$.  For $T$, the permutation action on the ten
two-subsets has character $1+\chi_V+\chi_W$ by direct counting
of invariant two-subsets.  This proves every row of (S507.7).
Thus all seven $g_i$ are monomial tests.  The tests with
$H=G$ give $a,b\ge0$, and the quartic row gives $e\ge0$.
Together with the stated $c,d\ge0$ premise this proves the
necessity of (S507.2).

For completeness, the next two tables give every distinct induced
linear-character multiplicity vector, together with a nonnegative
integer decomposition of its order functional into (S507.2).
The columns labeled $a,\ldots,e$ are the coefficients of that
functional, not values of the order variables.  Thus each last
entry can be verified by substituting the definitions of $g_i$.
\[
\begin{array}{r|rrrrrrr|l}
 i&a&b&c&d&x&y&e&\text{nonnegative decomposition}\\\hline
1 & 0 & 0 & 0 & 0 & 0 & 0 & 1 & e \\
2 & 0 & 0 & 0 & 0 & 0 & 0 & 2 & 2e \\
3 & 0 & 0 & 0 & 0 & 1 & 1 & 0 & g_1 \\
4 & 0 & 0 & 0 & 1 & 0 & 0 & 1 & d+e \\
5 & 0 & 0 & 0 & 1 & 0 & 1 & 1 & g_5 \\
6 & 0 & 0 & 0 & 1 & 1 & 1 & 1 & g_1+d+e \\
7 & 0 & 0 & 1 & 0 & 0 & 0 & 1 & c+e \\
8 & 0 & 0 & 1 & 0 & 1 & 0 & 1 & g_4 \\
9 & 0 & 0 & 1 & 0 & 1 & 1 & 1 & g_1+c+e \\
10 & 0 & 0 & 1 & 1 & 0 & 0 & 2 & c+d+2e \\
11 & 0 & 0 & 1 & 1 & 1 & 1 & 1 & g_1+c+d+e \\
12 & 0 & 0 & 1 & 1 & 1 & 1 & 2 & g_4+g_5 \\
13 & 0 & 0 & 1 & 1 & 2 & 2 & 2 & g_1+g_4+g_5 \\
14 & 0 & 0 & 2 & 2 & 2 & 2 & 4 & 2g_4+2g_5 \\
15 & 0 & 1 & 0 & 0 & 0 & 0 & 0 & b \\
16 & 0 & 1 & 0 & 0 & 1 & 0 & 0 & g_2 \\
17 & 0 & 1 & 0 & 1 & 0 & 0 & 0 & b+d \\
18 & 0 & 1 & 0 & 1 & 0 & 1 & 0 & g_7 \\
19 & 0 & 1 & 0 & 1 & 1 & 1 & 0 & g_1+b+d \\
20 & 0 & 1 & 0 & 2 & 0 & 1 & 1 & g_7+d+e \\
21 & 0 & 1 & 0 & 2 & 1 & 2 & 1 & g_1+g_7+d+e \\
22 & 0 & 1 & 1 & 1 & 0 & 1 & 1 & g_7+c+e \\
\end{array}
\]

\[
\begin{array}{r|rrrrrrr|l}
 i&a&b&c&d&x&y&e&\text{nonnegative decomposition}\\\hline
23 & 0 & 1 & 1 & 1 & 2 & 1 & 1 & g_1+g_4+b+d \\
24 & 0 & 1 & 1 & 3 & 2 & 3 & 3 & g_1+g_4+g_5+g_7+d+e \\
25 & 1 & 0 & 0 & 0 & 0 & 0 & 0 & a \\
26 & 1 & 0 & 0 & 0 & 0 & 1 & 0 & g_3 \\
27 & 1 & 0 & 1 & 0 & 0 & 0 & 0 & a+c \\
28 & 1 & 0 & 1 & 0 & 1 & 0 & 0 & g_6 \\
29 & 1 & 0 & 1 & 0 & 1 & 1 & 0 & g_1+a+c \\
30 & 1 & 0 & 1 & 1 & 1 & 0 & 1 & g_6+d+e \\
31 & 1 & 0 & 1 & 1 & 1 & 2 & 1 & g_1+g_5+a+c \\
32 & 1 & 0 & 2 & 0 & 1 & 0 & 1 & g_6+c+e \\
33 & 1 & 0 & 2 & 0 & 2 & 1 & 1 & g_1+g_6+c+e \\
34 & 1 & 0 & 3 & 1 & 3 & 2 & 3 & g_1+g_4+g_5+g_6+c+e \\
35 & 1 & 1 & 0 & 0 & 0 & 0 & 0 & a+b \\
36 & 1 & 1 & 0 & 0 & 1 & 1 & 0 & g_2+g_3 \\
37 & 1 & 1 & 0 & 0 & 1 & 1 & 2 & g_2+g_3+2e \\
38 & 1 & 1 & 1 & 1 & 0 & 0 & 0 & a+b+c+d \\
39 & 1 & 1 & 1 & 1 & 1 & 1 & 0 & g_6+g_7 \\
40 & 1 & 1 & 1 & 1 & 2 & 2 & 0 & g_1+g_6+g_7 \\
41 & 1 & 1 & 2 & 2 & 1 & 1 & 2 & g_6+g_7+c+d+2e \\
42 & 1 & 1 & 2 & 2 & 3 & 3 & 2 & g_1+g_4+g_5+g_6+g_7 \\
43 & 1 & 1 & 4 & 4 & 5 & 5 & 6 & g_1+3g_4+3g_5+g_6+g_7 \\
\end{array}
\]

Here is an exhaustive finite algorithm for these tables, with no
numerical optimization premise.  List all $120$ permutations of
five letters and their exact composition table.  Begin with the
trivial subgroup; repeatedly adjoin every individual permutation
to every subgroup found, compute the resulting generated subgroup,
and discard duplicates.  This finds every subgroup: given a
subgroup $H$, successively adjoin elements of $H$ until its
finite generating list has been exhausted.  Each such step appears
in this search.  It yields $156$ subgroups and $19$ conjugacy
classes; conjugate subgroups give the same induced patterns.
For each representative, retain a generating list.

The exponent of $S_5$ divides $60$, so every linear character of
one of its subgroups has values among the $60$th roots of unity.
Assign to each generator all phases $j\in\mathbb Z/60$ allowed
by its order.  Propagate those phases through its Cayley graph,
reject inconsistent assignments, and check
$j(hh')=j(h)+j(h')$ for all pairs.  This both finds every linear
character and verifies that each retained assignment is one.
Compute (S506.7) in the exact ring
$\mathbb Q[Z]/(\Phi_{60}(Z))$, using
$\overline{\lambda(h)}=Z^{-j(h)}$; reduce every sum modulo
$\Phi_{60}$ and divide by $|H|$.  The resulting multiplicities
are the nonnegative integers listed above.  As checks, the number
of characters found for each $H$ is $|H/[H,H]|$, and the
weighted dimension of every row is $120/|H|$.
Discarding repeated vectors gives exactly the $43$ displayed rows.
All steps use finite permutations, integer arithmetic and exact
cyclotomic reduction.

The full executable enumeration and its output are
\texttt{work/checks/s5\_hecke\_cone.py} and
\texttt{s5\_hecke\_cone.json}; their exploratory floating-point
optimizer is not used in this proof.  The separate exact verifier
\texttt{work/checks/s5\_cone\_decomposition\_verification.py}
and its JSON output reconstruct all $43$ rows as nonnegative
integer combinations of the twelve generators and verify the
sharpness witnesses below.  Its input SHA-256 is
\begin{center}\small\ttfamily
0b1ab9c1b47311477e3e113cced1e568a4c6fc417c8077875e87e7d8637ea917
\end{center}
The exhaustive permutation/character portion was re-executed
without its optimizer, and the displayed decompositions were
checked in integer arithmetic.  This is a transparent finite
computation with its algorithm and complete output supplied;
it is not a proof-assistant certification.

Every order vector satisfying (S507.2) pairs nonnegatively
with every displayed row.  Exhaustiveness therefore proves
sufficiency for (S507.1).  This establishes the exact cone.
\end{proof}

\begin{proof}[Proof of the pole bounds, sharpness and support condition]
The regular character of $G$ has multiplicities equal to
irreducible dimensions, whereas the character for the quadratic
field $E$ is $\mathbf1+\eta$.  All-prime Artin factorization
therefore gives (S507.3), including at ramified primes.
If $x=-v<0$, the inequalities $g_1,g_2,g_4,g_6\ge0$
give the first line of (S507.4).  Since $d,e,x+y\ge0$,
\[
 q=4(c+e)+4d+5(x+y)+2e\ge4(c+e)\ge4v,
 \qquad r=a+b+q\ge5v.
\]
If $c=0$, $g_4,g_6\ge0$ force $e,a\ge v$, so
$q\ge6e\ge6v$ and $r\ge a+b+q\ge8v$.
Twisting interchanges $a,b$, $c,d$ and $x,y$, giving every
symmetric assertion.  Since $g_1\ge0$, $x,y$ cannot both be
negative.  The other five irreducible factors are already
holomorphic at $z$ under the stated hypotheses.  A remaining
pole has integral order $v\ge1$ and $q\ge4v$; this proves both
the $q\le3$ holomorphy criterion and the floor bound.

The integer vector
\[
        (a,b,c,d,x,y,e)=(0,1,1,0,-1,1,0)
 \tag{S507.8}
\]
satisfies (S507.2): the seven $g_i$ equal
$(0,0,1,0,1,0,2)$ and the five required coordinates are
nonnegative.  It has a pole order $v=1$, $q=4$ and $r=5$.
Its sign twist provides the symmetric witness.  For the
$c=0$ branch, the vector
\[
        (a,b,c,d,x,y,e)=(1,1,0,0,-1,1,1)
 \tag{S507.9}
\]
has $g_i=(0,0,2,0,2,0,2)$, $q=6$ and $r=8$.
Multiplying these vectors by any positive integer gives every
pole order $v$.  By the exact cone, all $43$ monomial tests are
nonnegative at these witnesses.  This proves sharpness within
the specified positivity information.  It makes no claim that
these integer vectors arise from number fields or analytic
$L$-functions.

Finally $x<0$ forces $b>0$.  If $c>0$ the first branch of
(S507.6) holds.  If $c=0$, then $a,e>0$, which is exactly
its other branch.  Since $c\ge0$, these cases exhaust all
possibilities.  The sign twist proves the second inclusion.
Under $A_D$, (S506.4) confines the poles to the critical line.
Under $M_H$ and the Calegari family hypotheses, MT-S505 excludes
every real point of the open critical strip, including the
centre.  This leaves only the nonreal points described above.
\end{proof}

\paragraph{Input separation and the geometric gate.}
MT-S506 is a new reduction using classical character theory,
Artin induction and Hecke analyticity; only its last pole-location
assertion invokes $A_D$.  MT-S507 additionally uses the
four-dimensional holomorphy premise, supplied in the indicated
family by Calegari's theorem.  Its stronger exclusion of real
poles uses the separately proved, $M_H$-conditional MT-S505.
The finite cone does not establish noncoincidence, pole absence,
automorphy of the rank-five factors, or a new proof of GRH.

On the arithmetic base with the ramified primes removed, these
representations are the Artin motives of finite covers and their
direct summands.  Induction corresponds to finite pushforward
from a cover; (S507.1) records positivity of the zero orders
for all one-dimensional induced summands.  The witness (S507.8)
shows precisely why that information, even supplemented by the
rank-four holomorphy input, does not force the order divisor of
each irreducible summand to be effective.  The additional
all-prime induction identities are compatible with inertia
invariants at the removed primes.  Thus (S507.6), with the order
penalties (S507.4)--(S507.5), is a smaller explicit interface for
the Artin-positivity gate, rather than an identification of
zero-independent positive cohomology.  The attached run is
\texttt{GRH\_NIGHT/complexity}; the inherited interfaces are
MT-ART07 and MT-S501/MT-S504/MT-S505, while the cone certificate
and the deductions in MT-S506/MT-S507 are new in this run.
